                                              IMC 2023
                                     First Day, August 2, 2023
                                                   Solutions


Problem 1. Find all functions f : R → R that have a continuous second derivative and for which
the equality f (7x + 1) = 49f (x) holds for all x ∈ R.
                                   (proposed by Alex Avdiushenko, Neapolis University Paphos, Cyprus)

Hint:

   • The fixed point of 7x + 1 is −1/6.

   • Differentiating twice cancels out the coefficient 49.


Solution. Differentiating the equation twice, we get
                                                                                     
                              ′′              ′′           ′′          ′′       x−1
                            f (7x + 1) = f (x) or f (x) = f                               .               (1)
                                                                                 7

   Take an arbitrary x ∈ R, and construct a sequence by the recurrence
                                                                xk − 1
                                          x0 = x,    xk+1 =            .
                                                                   7
By (1), the values of f ′′ at all points of this sequence are equal. The limit of this sequence is − 16 ,
since xk+1 + 61 = 71 xk + 16 .
    Due to the continuity of f ′′ , the values of f ′′ at all points of this sequence are equal to f ′′ − 61 ,
                                                                                                            

which means that f ′′ (x) is a constant.
    Then f is an at most quadratic polynomial, f (x) = ax2 + bx + c. Substituting this expression
into the original equation, we get a system of equations, from which we find a = 36c, b = 12c, and
hence
                                            f (x) = c(6x + 1)2 .




                                                       1
Problem 2. Let A, B and C be n × n matrices with complex entries satisfying

                               A2 = B 2 = C 2     and B 3 = ABC + 2I.

Prove that A6 = I.
                                                        (proposed by Mike Daas, Universiteit Leiden)

Hint: Factorize B 3 − ABC.

Solution. Note that B 3 = A2 B, from which it follows that

                            A2 B − ABC = 2I =⇒ A(AB − BC) = 2I.

Similarly, using that B 3 = BC 2 , we find that

                            BC 2 − ABC = 2I =⇒ (BC − AB)C = 2I.

It follows that A is a left-inverse of (AB − BC)/2, whereas −C is a right inverse. Hence A = −C
and as such, it must hold that ABA = 2I − B 3 . It follows that ABA must commute with B, and so
it follows that (AB)2 = (BA)2 . Now we compute that

  (AB − BA)(AB + BA) = (AB)2 + AB 2 A − BA2 B − (BA)2 = (AB)2 + A4 − B 4 − (AB)2 = 0.

However, we noted before that the matrix AB − BC = AB + BA must be invertible. As such, it
must follow that AB = BA. We conclude that ABA = A2 B = B 3 and so it readily follows that
B 3 = I. Finally, A6 = B 6 = (B 3 )2 = I 2 = I, completing the proof.




                                                    2
Problem 3. Find all polynomials P in two variables with real coefficients satisfying the identity

                                 P (x, y)P (z, t) = P (xz − yt, xt + yz).

                                   (proposed by Giorgi Arabidze, Free University of Tbilisi, Georgia)

Hint: The polynomials (x+iy)n and (x−iy)m are trivial complex solutions. Suppose that P (x, y) =
(x + iy)n (x − iy)m Q(x, y), where Q(x, y) is divisible neither by x + iy nor x = iy and consider Q(x, y).

Solution. First we find all polynomials P (x, y) with complex coefficients which satisfies the condition
of the problem statement. The identically zero polynomial clearly satisfies the condition. Let consider
other polynomials.
    Let i2 = −1 and P (x, y) = (x + iy)n (x − iy)m Q(x, y), where n and m are non-negative integers
and Q(x, y) is a polynomial with complex coefficients such that it is not divisible neither by x + iy
nor by x − iy. By the problem statement we have Q(x, y)Q(z, t) = Q(xz − yt, xt + yz). Note that
z = t = 0 gives Q(x, y)Q(0, 0) = Q(0, 0). If Q(0, 0) ̸= 0, then Q(x, y) = 1 for all x and y. Thus
P (x, y) = (x + iy)n (x − iy)m . Now consider the case when Q(0, 0) = 0.
    Let x = iy and z = −it. We have Q(iy, y)Q(−it, t) = Q(0, 0) = 0 for all y and t. Since Q(x, y)
is not divisible by x − iy, Q(iy, y) is not identically zero and since Q(x, y) is not divisible by x + iy,
Q(−it, t) is not identically zero. Thus there exist y and t such that Q(iy, y) ̸= 0 and Q(−it, t) ̸= 0
which is impossible because Q(iy, y)Q(−it, t) = 0 for all y and t.
    Finally, P (x, y) polynomials with complex coefficients which satisfies the condition of the problem
statement are P (x, y) = 0 and P (x, y) = (x + iy)n (x − iy)m . It is clear that if n ̸= m, then
P (x, y) = (x + iy)n (x − iy)m cannot be polynomial with real coefficients. So we need to require
n = m, and for this case P (x, y) = (x + iy)n (x − iy)n = (x2 + y 2 )n .
    So, the answer of the problem is P (x, y) = 0 and P (x, y) = (x2 + y 2 )n where n is any non-negative
integer.




                                                    3
Problem 4. Let p be a prime number and let k be a positive integer. Suppose that the numbers
ai = ik + i for i = 0, 1, . . . , p − 1 form a complete residue system modulo p. What is the set of possible
remainders of a2 upon division by p?
                                             (proposed by Tigran Hakobyan, Yerevan State University, Armenia)
               p−1
               Y
Hint: Consider     (ik + i).
                        i=0

Solution. First observe that p = 2 does not satisfy the condtion, so p must be an odd prime.
Lemma. If p > 2 is a prime and Fp is the field containing p elements, then for any integer 1 ≤ n < p
one has the following equality in the field Fp

                                                         p−1
                                            
                            Y               0, if                 is even
                                       n             gcd(p − 1, n)
                                 (1 + α ) =
                                             n
                           α∈F∗p               2 , otherwise

Proof. We may safely assume that n|p − 1 since it can be easily proved that the set of n-th powers
of the elements of F∗p coincides with the set of gcd(p − 1, n)-th powers of the same elements. Assume
that t1 , t2 , ..., tn are the roots of the polynomial tn + 1 ∈ Fp [x] in some extension of the field Fp . It
follows that
               Y                         n
                                        YY             n Y
                                                       Y               n Y
                                                                       Y              n
                                                                                      Y
                      (1 + αn ) =          (α − ti ) =     (α − ti ) =     (ti − α) =   Φ(ti ),
              α∈F∗p                 α∈F∗p i=1                   i=1 α∈F∗p           i=1 α∈F∗p        i=1

                                                         p−1
                                    Q
where we define Φ(t) =                  α∈F∗p (t − α) = t    − 1. Therefore

                        n            n                 n
                                                                                                (
     Y
                   n
                        Y   p−1
                                     Y
                                          n p−1
                                                       Y      p−1                                0, if p−1
                                                                                                        n
                                                                                                           is even
             (1 + α ) =   (ti − 1) =   ((ti ) n − 1) =   ((−1) n − 1) =                           n
     α∈F∗p                    i=1                 i=1                       i=1
                                                                                                 2 , otherwise

    Let us now get back to our problem. Suppose the numbers ik + i, 0 ≤ i ≤ p − 1 form a complete
residue system modulo p. It follows that
                                       Y             Y
                                          (αk + α) =    α
                                                        α∈F∗p               α∈F∗p


so that α∈F∗p (αk−1 + 1) = 1 in Fp . According to the Lemma, this means that 2k−1 = 1 in Fp , or
        Q

equivalently, that 2k−1 ≡ 1(mod p). Therefore a2 = 2k + 2 ≡ 4(mod p) so that the remainder of a2
upon division by p is either 4 when p > 3 or is 1, when p = 3.




                                                                      4
Problem 5. Fix positive integers n and k such that 2 ≤ k ≤ n and a set M consisting of n fruits.
A permutation is a sequence x = (x1 , x2 , . . . , xn ) such that {x1 , . . . , xn } = M . Ivan prefers some (at
least one) of these permutations. He realized that for every preferred permutation x, there exist k
indices i1 < i2 < . . . < ik with the following property: for every 1 ≤ j < k, if he swaps xij and xij+1 ,
he obtains another preferred permutation.
    Prove that he prefers at least k! permutations.
                                   (proposed by Ivan Mitrofanov, École Normale Superieur Paris)
                                                                                P −1
Hint: For every permutation z of M , choose a preferred permutation x such that   x (m)z −1 (m)
                                                                                                  m∈M
is maximal.

Solution. Let S be the set of all n! permutations of M , and let P be the set of preferred permutations.
For every permutation x ∈ S and m ∈ M , let x−1 (m) denote the unique number i ∈ {1, 2, . . . , n}
with xi = m.
   For every x ∈ P , define
                                          X                       X                    
                                                 −1       −1             −1       −1
             A(x) = z ∈ S : ∀y ∈ P              x (m)z (m) ≥            y (m)z (m) .
                                                 m∈M                        m∈M


                                                                                                       x−1 (m)z −1 (m) is
                                                                                                 P
For every permutation z ∈ S, we can choose a permutation x ∈ P for which
                                                                                                 m∈M
maximal, and then we have z ∈ A(x); hence, all z ∈ S is contained in at least one set A(x).
                                               n!
   So, it suffices to prove that A(x) ≤           for every preferred permutation x. Fix x ∈ P , and
                                               k!
consider an arbitrary z ∈ A(x). Let the indices i1 < . . . < ik be as in the statement of the problem,
and let mj = xij for j = 1, 2, . . . , k.
   For s = 1, 2, . . . , k − 1 consider the permutation y obtained from x by swapping ms and ms+1 .
Since y ∈ P , the definition of A(x) provides

                        is z −1 (ms ) + is+1 z −1 (ms+1 ) ≥ is+1 z −1 (ms ) + is z −1 (ms+1 ),

                                             z −1 (ms+1 ) ≥ z −1 (ms ).
Therefore, the elements m1 , m2 , . . . , mk appear in z in this order. There are exactly n!/k! permuta-
                                           n!
tions with this property, so A(x) ≤ .
                                           k!




                                                          5
